\section{Consequences and open problems}\label{sec:consequences}

\subsection{Conjecture 15 of~\cite{DEFH2023}}

In a strong I-contractible graph the order of deletions matters~\cite[Fig.~3]{DEFH2023}; Theorem~16 of~\cite{DEFH2023} shows that for vertex-strong I-contractible graphs a bad choice of an $\IVS$-deletable vertex first occurs at $12$ vertices. For edges, Conjecture~15 of~\cite{DEFH2023} proposed that the first bad choice of an $\IS$-deletable edge occurs in the $15$-vertex graph of their Fig.~2B (the dunce hat $\Dunce$ plus one edge).

\begin{theorem}\label{thm:D}
The smallest graph $G$ (in terms of vertices) such that (1) $G\in\IS$ and (2) $G$ has an edge $\{v,w\}$ with $\Nb(v,w)\in\IS$ and $G-\{v,w\}\notin\IS$ has $14$ vertices. The graphs $\Ga+\{0,1\}$, $\Gb+\{0,3\}$ and $\Gc+\{2,3\}$ (\gsix{M\_b@bFHNEk@\textasciitilde{}|k\}U?}, \gsix{MCB@thku?\textasciitilde{}Lylk\}[?}, \gsix{M@r@d`ifbiFfo\}\textasciitilde{}o?}) are such graphs, and there are $32$ pairwise non-isomorphic such graphs in the $\mathrm{I}_4$-closure of $\Ga$ alone.
\end{theorem}
\begin{proof}
$\Ga+\{0,1\}\in\IS$ with $\Nb(0,1)\in\IS$ and $(\Ga+\{0,1\})-\{0,1\}=\Ga\notin\IS$ by Theorem~\ref{thm:B}; likewise for the other two. For minimality, if $G$ on $n\le 13$ vertices satisfied (1) and (2), then $G-\{v,w\}$ would be $\Ztwo$-acyclic (the deletion is a homotopy equivalence since $\cl(\Nb(v,w))$ is contractible) and not in $\IS$, so $n=13$ and $G-\{v,w\}$ is one of the three graphs of Theorem~\ref{thm:A}, by Lemma~\ref{lem:stuck}(a) a graph in the $\mathrm{I}_4$-closure of $\Gt$; but $G$ itself would then be in that closure, which contains no member of $\IS$. The count of $32$ is by direct computation of the closure of $\Ga$ (\texttt{data/}).
\end{proof}

Thus ``the order matters'' for $\IS$-deletable edges first at $14$ vertices, two more than for $\IVS$-deletable vertices.

\subsection{The conjectures of~\cite{FriasArmenta2020}}

Fr\'ias-Armenta~\cite{FriasArmenta2020} gave a $13$-vertex counterexample to an axiom used by Ivashchenko and isolated three statements of~\cite{Ivashchenko1994a} whose published proofs relied on it, restating them as conjectures: (Conjecture~3) every $G\in\I$ on $n$ vertices can be transformed into the complete graph $K_n$ by additions of edges $uv$ with $\Nb(u,v)\in\I$; (Conjecture~4) if $G\in\I$ and $G-v\in\I$ then $\Nb(v)\in\I$; (Conjecture~5) if $G\in\I$ then $\Nb(v)$ and $G-v$ are I-homotopy equivalent. Our examples are natural test cases, being the smallest graphs where $\I$ and $\IS$ differ, but they do not refute these conjectures.

\begin{proposition}\label{prop:frias}
Each of $\Ga$, $\Gb$, $\Gc$ and $\Dunce$ can be transformed into a cone, hence into $K_n$, by $11$, $12$, $12$ and $8$ additions of edges $uv$ with $\Nb(u,v)\in\IS$ respectively (sequences in \ref{app:certificates}). For every vertex $v$ of these four graphs, $\widetilde H_*(G-v;\ZZ)\cong\widetilde H_*(\Nb(v);\ZZ)$, and no $G-v$ is contractible.
\end{proposition}

The second statement is not a coincidence: if $\cl(G)$ is contractible then $\cl(G)=\cl(G-v)\cup\mathrm{cone}(\cl(\Nb(v)))$ with intersection $\cl(\Nb(v))$, and the Mayer--Vietoris sequence gives $\widetilde H_k(G-v)\cong\widetilde H_k(\Nb(v))$ for all $k$. Hence Conjectures~4 and~5 cannot be refuted by homological invariants on graphs in $\I$; a counterexample would have to be detected by the fundamental group or by simple homotopy type. (By~\cite{DochtermannMatsushita2025}, $\I$-homotopy equivalence of graphs is simple homotopy equivalence of clique complexes, so Conjecture~5 asks whether $\cl(G-v)$ and $\cl(\Nb(v))$ are simple homotopy equivalent whenever $\cl(G)$ is contractible.)

\subsection{Open problems}\label{sec:open}

\begin{question}\label{q:K4}
Is every graph in $\I\setminus\IS$ on $14$ vertices $K_4$-free? Equivalently, is the list of eight graphs of Section~\ref{sec:fourteen} complete?
\end{question}

The full class of $2$-connected graphs on $14$ vertices with $4\le\delta\le\Delta\le 12$ has about $10^{16}$ members, two to three orders of magnitude beyond the present computation. Two routes seem feasible: enumerating only graphs with a bounded number of copies of $K_4$, or growing the stuck graphs on $13$ and $14$ vertices by one vertex and testing stuckness of the extensions.

\begin{conjecture}\label{conj:disc}
Every stuck graph whose clique complex is $2$-dimensional unfolds, after cutting along the edges in three triangles, to a disc whose boundary is identified with an induced cycle $C$ of $G$ with the pattern $a\,a\,a^{-1}$ (contractible case) or $a^{3}$ (torsion case), possibly with pendant edges traversed back and forth.
\end{conjecture}

Among the $51$ stuck $K_4$-free graphs on $14$ vertices, $48$ satisfy the conjecture; the three exceptions (two annuli and one disconnected unfolding) all have $45$ or $47$ edges and torsion, and it may be that the conjecture holds for stuck graphs of minimum size.

\begin{question}
Are the clique complexes of $\Ga$, $\Gb$, $\Gc$ collapsible? Conjecture~27 of~\cite{EFH2022} asserts that a flag complex is collapsible if and only if its $1$-skeleton is in $\IS$; since the three graphs are not in $\IS$, either they are non-collapsible (as the dunce hat is) or they refute that conjecture.
\end{question}

\begin{question}
Theorem~\ref{thm:A} and Adamaszek's results~\cite{Adamaszek2014} give $13$ as the smallest number of vertices of a flag complex with odd torsion in $H_1$ \emph{among $\Ztwo$-acyclic flag complexes}. Is $13$ also the minimum without the $\Ztwo$-acyclicity assumption, and which torsion groups occur on $13$ and $14$ vertices?
\end{question}

\begin{question}
Table~\ref{tab:card} shows $\IVS=2\text{-Dis}=3\text{-Dis}$ for $n\le 12$. Are there graphs in $\IVS$ that are not $2$-dismantlable, and what is the smallest one?
\end{question}

Finally, Question~28 of~\cite{DEFH2023} asked whether every graph with contractible clique complex is I-contractible; this is now the theorem of Dochtermann and Matsushita~\cite{DochtermannMatsushita2025}. Our results show that the analogous question with ``$\Ztwo$-acyclic'' in place of ``contractible'' has a negative answer already at $13$ vertices, and that the gap between $\I$ and $\IS$ opens at $14$.
